% ======================================================================
% Diagrama de flujo — generado automáticamente por generar_diagrama_flujo.py
% NO EDITAR MANUALMENTE; regenerar desde el descriptor JSON.
% ======================================================================
\documentclass[11pt,a4paper]{article}

% ── Codificación, fuentes y lenguaje ───────────────────────────────────────────
\usepackage[utf8]{inputenc}
\usepackage[T1]{fontenc}
\usepackage[default,scale=0.95]{sourcesanspro}
\renewcommand{\familydefault}{\sfdefault}
\usepackage[spanish,es-noshorthands,es-tabla]{babel}

% ── Geometría y espaciado ─────────────────────────────────────────────────────
\usepackage[top=2.2cm, bottom=2.5cm, left=2.2cm, right=2.2cm, headheight=15pt]{geometry}
\usepackage{setspace}
\setstretch{1.18}
\usepackage{parskip}

% ── TikZ (simbolos ISO 5807 / ANSI X3.5) ──────────────────────────────────────
\usepackage{tikz}
\usetikzlibrary{babel, arrows.meta, positioning, shapes.geometric,
                shapes.symbols, decorations.pathmorphing, calc}

% ── Tablas y listas ───────────────────────────────────────────────────────────
\usepackage{booktabs}
\usepackage{tabularx}
\usepackage{array}
\usepackage{enumitem}

% ── Colores, cajas e iconos ───────────────────────────────────────────────────
\usepackage[dvipsnames]{xcolor}
\usepackage{tcolorbox}
\tcbuselibrary{skins, breakable}
\usepackage{fontawesome5}

% ── Secciones con barra de color (infografía) ─────────────────────────────────
\usepackage{titlesec}
\titleformat{\section}
  {\normalfont\LARGE\bfseries\color{azulNoche}}
  {\colorbox{cyanNeon}{\color{fondoOscuro}\thesection}}{0.7em}{}
  [\vspace{2pt}\color{cyanNeon}\rule{\linewidth}{1.8pt}]
\titlespacing*{\section}{0pt}{24pt}{10pt}
\titleformat{\subsection}
  {\normalfont\large\bfseries\color{azulNoche}}
  {\textcolor{cyanNeon}{\thesubsection}}{0.7em}{}
  [\vspace{1pt}\color{grisLinea}\rule{\linewidth}{0.6pt}]
\titlespacing*{\subsection}{0pt}{14pt}{6pt}

% ── Cabeceras y pies de página ────────────────────────────────────────────────
\usepackage{fancyhdr}
\usepackage{lastpage}
\pagestyle{fancy}
\fancyhf{}
\renewcommand{\headrulewidth}{0pt}
\renewcommand{\footrulewidth}{0pt}
\fancyfoot[C]{%
  \begin{minipage}{\linewidth}
    \vspace{2pt}
    \color{cyanNeon}\rule{\linewidth}{1.2pt}\\[2pt]
    \centering\color{grisTexto}\footnotesize
    Página \thepage\ de \pageref{LastPage}
  \end{minipage}%
}

% ── Hiperenlaces ──────────────────────────────────────────────────────────────
\usepackage[colorlinks=true, linkcolor=azulNoche, urlcolor=cyanNeon]{hyperref}
\sloppy
\emergencystretch 3em

% ── Paleta de colores ─────────────────────────────────────────────────────────
\definecolor{fondoOscuro}{HTML}{0D1B2A}
\definecolor{verdeTurquesa}{HTML}{004D40}
\definecolor{azulNoche}{HTML}{1B3A6B}
\definecolor{azulMedio}{HTML}{2A5298}
\definecolor{cyanNeon}{HTML}{00C8E0}
\definecolor{verdeSignal}{HTML}{00B887}
\definecolor{naranjaVivo}{HTML}{FF7043}
\definecolor{rojoAlerta}{HTML}{E53935}
\definecolor{grisPapel}{HTML}{F4F6FA}
\definecolor{grisLinea}{HTML}{C8D0E0}
\definecolor{grisTexto}{HTML}{6B7A99}
\definecolor{amarilloNota}{HTML}{FFB300}

% ── Tarjetas ──────────────────────────────────────────────────────────────────
\tcbset{
  cajaContenido/.style={
    enhanced, breakable,
    arc=5pt, outer arc=5pt,
    colback=grisPapel, colframe=grisLinea,
    boxrule=0.8pt,
    fonttitle=\bfseries\small, coltitle=azulNoche,
    top=6pt, bottom=6pt, left=10pt, right=10pt,
  },
}

% ── Macros ────────────────────────────────────────────────────────────────────
\newcommand{\iconotexto}[2]{\textcolor{cyanNeon}{\faIcon{#1}}~#2}
\newcommand{\iconoBanda}{sitemap}
\newcommand{\bandaTitulo}[3][fondoOscuro]{%
  \noindent
  \begin{tcolorbox}[%
    enhanced, arc=4mm, colback=#1!10!white, colframe=#1, boxrule=0pt,
    borderline west={6pt}{0pt}{cyanNeon},
    top=6pt, bottom=6pt, left=8pt, right=8pt,
  ]
    \noindent
    \begin{minipage}[c]{0.10\linewidth}
      \centering\color{cyanNeon}\faIcon{sync-alt}
    \end{minipage}%
    \hspace{8pt}%
    \begin{minipage}[c]{0.88\linewidth}
      {\fontsize{20}{24}\selectfont\bfseries\color{white}#2}\par\vspace{5pt}%
      \textcolor{cyanNeon!80!white}{\small\faIcon{angle-right}~#3}%
    \end{minipage}
    \vspace{4pt}
  \end{tcolorbox}%
  \vspace{8pt}%
}


  % ── Estilos de flujo (ISO 5807 / ANSI X3.5) ─────────────────────────────
  \tikzset{
    flujo/.style={
      >=Stealth,
      font=\footnotesize\sffamily,
    },
    terminador/.style={
      draw=azulNoche, fill=azulNoche!8!white, rounded corners=3pt,
      text width=2.2cm, align=center, inner sep=2mm,
      font=\footnotesize\sffamily,
    },
    proceso/.style={
      draw=azulNoche, fill=cyanNeon!10!white,
      text width=2.2cm, align=center, inner sep=2mm,
      font=\footnotesize\sffamily,
    },
    entradasalida/.style={
      draw=azulNoche, fill=verdeSignal!10!white,
      trapezium, trapezium left angle=75, trapezium right angle=105,
      text width=2.2cm, align=center, inner sep=2mm,
      font=\footnotesize\sffamily,
    },
    decision/.style={
      draw=azulNoche, fill=amarilloNota!20!white,
      diamond, aspect=1.6, text width=2.0cm, align=center, inner sep=1pt,
      font=\footnotesize\sffamily,
    },
    almacenamiento/.style={
      draw=azulNoche, fill=grisPapel,
      cylinder, shape border rotate=90, aspect=0.3,
      text width=2.2cm, align=center, inner sep=2mm,
      font=\footnotesize\sffamily,
    },
    documento/.style={
      draw=azulNoche, fill=azulMedio!8!white,
      text width=2.2cm, align=center, inner sep=2mm,
      font=\footnotesize\sffamily,
    },
    nota/.style={
      draw=grisLinea, dashed, fill=grisPapel,
      text width=2.4cm, align=center, inner sep=2mm,
      font=\scriptsize\sffamily,
    },
    etiqueta/.style={
      font=\scriptsize\sffamily, fill=white,
      fill opacity=0.75, text opacity=1, inner sep=1pt,
    },
  }


% ── Cabecera del documento ────────────────────────────────────────────────────
\fancyhead[L]{\small\color{azulNoche}\textbf{ELECTRÓNICA}}
\fancyhead[R]{\small\color{grisTexto}Sincronización FAQ → Flujo}
\renewcommand{\headrulewidth}{0pt}

% ─────────────────────────────────────────────────────────────────────────────
\begin{document}

\bandaTitulo{Sincronización FAQ → Flujo}{YAML, hash determinista y re-traducción quirúrgica de un documento único de higiene de estado}

\vspace{4pt}
\begin{center}
\begin{tcolorbox}[cajaContenido, title={\faIcon{map-signs}~Leyenda de símbolos---ISO 5807 / ANSI X3.5}]
\small
Óvalo = \textbf{terminador} (inicio/fin) $\cdot$
Rectángulo = \textbf{proceso} (acción determinista) $\cdot$
Rombo = \textbf{decisión} (ramas Sí/No) $\cdot$
Paralelogramo = \textbf{entrada/salida} (lectura/escritura de archivos) $\cdot$
Cilindro = \textbf{almacenamiento online} (estado, snapshots) $\cdot$
Rectángulo con base ondulada = \textbf{documento} (sesiones, logs) $\cdot$
Rectángulo punteado gris = \textbf{nota explicativa} (sin acción).
Flujo: arriba$\rightarrow$abajo, izquierda$\rightarrow$derecha.

\faIcon{tags}~Cada símbolo lleva una \textbf{etiqueta} (T/P/D/E/A/N + número, según su tipo);
su significado se expande en la tabla \emph{Leyenda de etiquetas} bajo cada diagrama.
\end{tcolorbox}
\end{center}

\vspace{6pt}

\section{\iconotexto{sync-alt}{Fase única: actualización faq\_higiene\_estado\_sesion.md}}

\begin{center}
\begin{tikzpicture}[flujo, >=Stealth, x=1cm, y=1cm]
  % Fondo decorativo
  \fill[azulNoche!3!white, rounded corners=4pt]
    (-2.3,1.8) rectangle (12.8,-18.3);
  \node[terminador] (T1) at (0.0,0.0) {Inicio};
  \node[proceso] (P1) at (0.0,-2.0) {Verificar
daías y
secciones};
  \node[decision] (D1) at (0.0,-4.0) {adelanto =
'detalle'?};
  \node[almacenamiento] (A1) at (3.5,-4.0) {extracto.md
nuevo};
  \node[proceso] (P2) at (3.5,-6.0) {Calcular
hash sha256};
  \node[almacenamiento] (A2) at (0.0,-6.0) {extracto.md
anterior};
  \node[decision] (D2) at (0.0,-8.0) {hash igual?};
  \node[entradasalida] (E1) at (3.5,-8.0) {Sincronizar
YAML};
  \node[decision] (D3) at (3.5,-10.0) {cambio
estructural?
LLM};
  \node[terminador] (T2) at (0.0,-10.0) {Fin
(proceso inerte)};
  \node[proceso] (P3) at (7.0,-10.0) {Re-traducir
de pasajes
semánticos};
  \node[proceso] (P4) at (7.0,-12.0) {Traducir
quiúrgico
al flujo};
  \node[proceso] (P5) at (10.5,-12.0) {Abortar
con aviso};
  \node[almacenamiento] (A3) at (3.5,-14.0) {flujo.tex +
flujo.pdf};
  \node[terminador] (T3) at (7.0,-16.0) {Fin
(flujo
actualizado)};

  \draw[->] (T1.south) -- (P1.north);
  \draw[->] (P1.south) -- (D1.north);
  \draw[->] (D1.east) -- (A1.west) node[etiqueta, above]{adv:
'detalle'};
  \draw[->] (D1.south) -- (A2.north) node[etiqueta, left]{extracto
anterior};
  \draw[->] (A1.south) -- (P2.north) node[etiqueta, left]{nuevo};
  \draw[->] (A2.south) -- (D2.north);
  \draw[->] (P2.west) -- (D2.east);
  \draw[->] (D2.south) -- (T2.north) node[etiqueta, left]{sí
(no cambia)};
  \draw[->] (D2.east) -- (E1.west) node[etiqueta, above]{no
(cambió)};
  \draw[->] (E1.south) -- (D3.north);
  \draw[->] (D3.east) -- (P4.west) node[etiqueta, above]{estructural +
nuevo};
  \draw[->] (D3.east) -- (P3.west) node[etiqueta, above]{semántico};
  \draw[->] (P3.south) -- (P4.north);
  \draw[->] (P4.west) -- (A3.east);
  \draw[->] (D3.east) -- (P5.west) node[etiqueta, above]{--no-llm};
  \draw[->] (P5.west) -- (T3.east);
\end{tikzpicture}
\end{center}

\begin{tcolorbox}[cajaContenido, title={\faIcon{tags}~Leyenda de etiquetas}]
\small
\begin{tabularx}{\linewidth}{@{}llX@{}}
\toprule
\textbf{Etiqueta} & \textbf{Símbolo} & \textbf{Significado} \\
\midrule
A1 & Cilindro (almacenamiento) & snapshot del documento en estado detectado \\
A2 & Cilindro (almacenamiento) & version del documento de la iteracion previa \\
A3 & Cilindro (almacenamiento) & Entregables LaTeX del flujo, regenerados \\
D1 & Rombo (decisión) & Primera invocación del fujo (adelanto = 'detalle') \\
D2 & Rombo (decisión) & Compara hash anterior y nuevo del snapshot \\
D3 & Rombo (decisión) & Clasificacion por LLM de la diferencia (cosmetica vs estructural) \\
E1 & Paralelogramo (entrada/salida) & Reescritura del YAML maestro si el hash cambio \\
P1 & Rectángulo (proceso) & Revisión de higiene de sesión (bitácoras y estado) \\
P2 & Rectángulo (proceso) & Funcion pura y determinista sobre el contenido \\
P3 & Rectángulo (proceso) & Se aplica la re-traduccion LLM a pasajes semanticos \\
P4 & Rectángulo (proceso) & Propagacion del cambio estructural al flujo \\
P5 & Rectángulo (proceso) & Con --no-llm aborta; sin bandera, se detecta y documenta \\
T1 & Óvalo (terminador) & Punto de entrada del flujo \\
T2 & Óvalo (terminador) & Sin cambios; el documento y el flujo quedan sincronizados \\
T3 & Óvalo (terminador) & Entrega el flujo actualizado y compilado \\
\bottomrule
\end{tabularx}
\end{tcolorbox}


\section{\iconotexto{sticky-note}{Notas}}
\begin{itemize}
\item El flujo es 100\% determinista: proceso inerte si no hay cambios.
\item La re-traducción LLM solo aplica cambios estructurales aprobados; el modo --no-llm aborta el proceso ante cambios estructurales.
\item El diagrama de flujo es generado por execution/generar\_diagrama\_flujo.py a partir de este descriptor JSON.
\end{itemize}

% ─────────────────────────────────────────────────────────────────────────────
\end{document}
